\documentclass{article}
\usepackage[T1]{fontenc}
\usepackage{lmodern}
\usepackage{graphicx}
\usepackage{amsmath,amssymb}
\usepackage[a4paper,margin=2cm]{geometry}
\usepackage[absolute,overlay]{textpos}
\setlength{\TPHorizModule}{1mm}
\setlength{\TPVertModule}{1mm}
\usepackage[indonesian]{babel}
\usepackage{hyperref}
\setlength{\emergencystretch}{2em}
% unit: Halaman 37

\begin{document}

% === Halaman 37 (llm) ===

\section*{Elliptic Curve Cryptography (ECC) $\ast$)}

\begin{itemize}
    \item ECC adalah sistem kriptografi kunci-publik, sejenis dengan RSA, Rabin, ElGamal, D-H, dll.
    \item Setiap pengguna memiliki \textbf{kunci publik dan kunci privat}
    \begin{itemize}
        \item Kunci publik untuk enkripsi atau untuk verifikasi tanda tangan digital
        \item Kunci privat untuk dekripsi atau untuk menghasilkan tanda tangan digital
    \end{itemize}
    \item Kurva eliptik digunakan sebagai perluasan sistem kriptografi kunci-publik yang lain:
    \begin{enumerate}
        \item Elliptic Curve Elgamal (ECEG)
        \item Elliptic Curve Digital Signature (ECDSA)
        \item Elliptic Curve Diffie-Hellman (ECDH)
    \end{enumerate}
\end{itemize}

\vspace{10mm}

\textcolor{red}{\ast}) Sumber bahan: \textcolor{red}{Debdeep Mukhopadhyay, \textit{Elliptic Curve Cryptography}, Dept of Computer Sc and Engg IIT Madras}

\end{document}
